\section*{Capítulo I, Problema XVI}

\textbf{Problema:}

16. Sean $A, B$ vectores no nulos, mutuamente perpendiculares. Demuestra que para cualquier número $c$ se cumple que $\|A + cB\| \geq \|A\|$.

\textbf{Solución:}

Dado que $A$ y $B$ son mutuamente perpendiculares, se cumple que $A \cdot B = 0$. Por lo tanto, podemos calcular la norma de $A + cB$ como sigue:
\begin{align*}
    \|A + cB\|^2 &= (A + cB) \cdot (A + cB) \\
    &= A \cdot A + 2c(A \cdot B) + c^2(B \cdot B).
\end{align*}

Dado que $A \cdot B = 0$, el término $2c(A \cdot B)$ se anula. Por lo tanto:
\begin{align*}
    \|A + cB\|^2 &= \|A\|^2 + c^2\|B\|^2.
\end{align*}

Como $c^2\|B\|^2 \geq 0$, se tiene que:
\begin{align*}
    \|A + cB\|^2 \geq \|A\|^2.
\end{align*}

Tomando la raíz cuadrada en ambos lados, obtenemos:
\begin{align*}
    \|A + cB\| \geq \|A\|.
\end{align*}

Por lo tanto, se cumple la desigualdad $\|A + cB\| \geq \|A\|$ para cualquier número $c$.